\clearpage\fancyhead[L]{\small COMET.}\hypertarget{native-340-lower}{}\label{native-340-lower}
\begin{center}\textbf{COMET.}\par\vspace{4pt}\hrule\vspace{4pt}\fontsize{10}{12}\selectfont Monatliche Correspondenz. Band XXIV. S. 304\,..\,307. 1811 September.\par\vspace{4pt}\hrule\vspace{4pt}\end{center}

\par\indent ---\;---\;--- Interessante Resultate der Beobachtung und der Rechnung theilte uns Herr Professor \textsc{Gauss} mit. Aus einigen seiner Briefe heben wir das hieher gehörige aus: ---\;---\;--- Göttingen am 8. Sept. „Den Cometen, welchen ich zuerst am 22. August gesehen habe, konnte ich, da die Aussicht auf der Sternwarte nach Norden durch die Stadt sehr beschränkt ist, erst am 3. September zum erstenmal beobachten. Ich habe versucht, seinen Ort durch Abstände mit dem Sextanten zu bestimmen, welches mit dem schönen Stativ, was sie uns durch den geschickten Mechanicus \textsc{Körner} in Weimar besorgt haben, eine sehr bequeme und wie es scheint auch verhältnissmässig sehr genaue Beobachtungsart ist. Die Beobachtungen vom 3. Sept. habe ich noch nicht reducirt, aber die Beobachtungen der folgenden Tage, wobei zur Vergleichung, immer $\alpha$ Aurigae und $\alpha$ Lyrae genommen wurden, habe ich mit aller Sorgfalt berechnet und folgende Resultate gefunden:
\Needspace{9\baselineskip}
\begin{center}\fontsize{10}{12}\selectfont\setlength{\tabcolsep}{4pt}\renewcommand{\arraystretch}{1.15}\begin{tabular}{r|c|c|c}
1811 & \shortstack{Mittlere Zeit\\in Göttingen} & AR. Comet. & \shortstack{Nördliche\\Abweichung}\\\hline
Septbr. 4 & $8^{\mathrm u}\,28^{\mathrm m}\,47^{\mathrm s}$ & $158^\circ\,25'\,24''$ & $39^\circ\,18'\,2''$\\
6 & $8\;48\;38$ & $160\;23\;16$ & $40\;14\;16$\\
7 & $8\;57\;6$ & $161\;26\;10$ & $40\;41\;54$\\
\end{tabular}\end{center}

\SourceLeaf{341}\par\indent Ich habe nach diesen und neun Beobachtungen des Herrn von \textsc{Zach} meine parabolischen Elemente\footnote{Monatl. Corresp. 1811 August-Heft S. 180. [S. 326 d. B.]} verbessert, und so die folgenden herausgebracht:
\Needspace{10\baselineskip}
\begin{center}\fontsize{10}{12}\selectfont\setlength{\tabcolsep}{4pt}\renewcommand{\arraystretch}{1.15}\begin{tabular}{l}
Zeit der Sonnennähe 1811 Sept. 12\\
\quad $5^{\mathrm u}\,21^{\mathrm m}\,15^{\mathrm s}$ Mittlere Zeit in Göttingen\\
\end{tabular}\end{center}
\begin{center}\fontsize{10}{12}\selectfont\setlength{\tabcolsep}{4pt}\renewcommand{\arraystretch}{1.15}\begin{tabular}{llr}
Logarithm. des kleinsten Abstandes & .\;.\;.\;.\;.\;.\;.\;. & $0.017060$\\
Länge der Sonnennähe & .\;.\;.\;.\;.\;.\;.\;.\;.\;.\;.\;.\;. & $75^\circ\,17'\,34''$\\
Länge des aufsteigenden Knoten & .\;.\;.\;.\;.\;.\;.\;.\;. & $140\;24\;13$\\
Neigung der Bahn & .\;.\;.\;.\;.\;.\;.\;.\;.\;.\;.\;.\;.\;. & $73\;7\;16$\\
\end{tabular}\end{center}

\par\indent Von Ellipticität der Bahn scheint jetzt noch keine sichere Spur zu sein. Ich habe mir bisher nur ein paarmal die Zeit genommen, den Cometen mit einem grossen Telescope zu betrachten; am 7. Sept. war die Form des Schweifes sehr merkwürdig. Er bog sich in zwei Aeste vom Cometen ab, aber diese beiden Aeste gingen nicht vom Cometen selbst aus, sondern hingen in einer kleinen Entfernung von diesem, durch einen dunkeln Zwischenraum getrennt, zusammen, so dass sie den Comet wie eine Parabel ihren Brennpunkt einschlossen. Gewiss ist dieser Comet in mehr als einer Rücksicht einer der merkwürdigsten, die jemals beobachtet sind. Von einem eigentlichen Kern im Cometen, konnte ich eben so wenig, als Herr Prof. \textsc{Harding} eine Spur wahrnehmen.“ ---\quad Göttingen am 16. Sept. Ich werde nun, da Sie den Cometen im Meridian beobachten, meine Beobachtungen heute Abend schliessen. Hier was ich seit meinem letzten Brief erhielt:
\Needspace{10\baselineskip}
\begin{center}\fontsize{10}{12}\selectfont\setlength{\tabcolsep}{4pt}\renewcommand{\arraystretch}{1.15}\begin{tabular}{r|c|c|c}
 & \shortstack{Mittlere Zeit\\in Göttingen} & AR. Cometae & \shortstack{Nördliche\\Declination}\\\hline
Sept. 9 & $8^{\mathrm u}\,37^{\mathrm m}\,53^{\mathrm s}$ & $163^\circ\,35'\,5''$ & $41^\circ\,39'\,13''$\\
14 & $8\;45\;5$ & $169\;44\;58$ & .\quad.\quad.\quad.\\
15 & $7\;39\;57$ & $171\;3\;58$ & $44\;24\;57$\\
16 & $7\;57\;16$ & $172\;31\;4$ & $44\;51\;23$\\
\end{tabular}\end{center}

\Needspace{24\baselineskip}\par\indent Für den Lauf des Cometen bis Ende December habe ich nach meinen verbesserten Elementen folgende Ephemeride berechnet:
\begin{center}\fontsize{10}{12}\selectfont\setlength{\tabcolsep}{3pt}\renewcommand{\arraystretch}{1.15}\begin{tabular}{rr|c|c|c|c}
\multicolumn{2}{c|}{\shortstack{Mittlere Zeit\\in Göttingen}} & \shortstack{Gerade\\Aufsteigung} & \shortstack{Nördliche\\Abweichung} & \shortstack{Log. d. Entf.\\des Cometen\\von der Erde} & \shortstack{Licht-\\stärke}\\\hline
1811 Sept. 12 & $5^{\mathrm u}$ & $166^\circ\,54'$ & $42^\circ\,58'$ & $0.2079$ & $0.355$\\
17 & $0$ & $173\;46$ & $45\;15$ & $0.1836$ & $0.394$\\
22 & $6$ & $182\;2$ & $47\;17$ & $0.1594$ & $0.432$\\
.27 & $6$ & $192\;19$ & $48\;48$ & $0.1366$ & $0.465$\\
Oct. 2 & $7$ & $203\;24$ & $49\;34$ & $0.1160$ & $0.490$\\
7 & $7$ & $215\;58$ & $49\;7$ & $0.1004$ & $0.501$\\
12 & $8$ & $228\;45$ & $47\;16$ & $0.0911$ & $0.492$\\
22 & $9$ & $251\;23$ & $39\;46$ & $0.0963$ & $0.421$\\
Nov. 1 & $10$ & $267\;58$ & $30\;32$ & $0.1312$ & $0.310$\\
11 & $11$ & $279\;38$ & $22\;18$ & $0.1851$ & $0.208$\\
21 & $12$ & $288\;13$ & $14\;55$ & $0.2415$ & $0.138$\\
Dec. 1 & $13$ & $294\;54$ & $9\;42$ & $0.2976$ & $0.093$\\
11 & $14$ & $300\;24$ & $5\;49$ & $0.3487$ & $0.064$\\
21 & $15$ & $305\;7$ & $2\;55$ & $0.3938$ & $0.046$\\
31 & $16$ & $309\;16$ & $0\;48$ & $0.4330$ & $0.034$\\
\end{tabular}\end{center}

\Needspace{16\baselineskip}\par\indent Eine ganz flüchtige Vergleichung meiner Beobachtungen und der ersten \textsc{Olbers}schen mit meinen Elementen, hat mir folgende Resultate gegeben;
\SourceLeaf{342}
\begin{center}\fontsize{10}{12}\selectfont\setlength{\tabcolsep}{7pt}\renewcommand{\arraystretch}{1.15}\begin{tabular}{lr|r|r}
 & & AR. & Decl.\\\hline
Aug. & 23 & $+127''$ & $+18''$\\
Sept. & 4 & $+30$ & $-9$\\
 & 6 & $+9$ & $-11$\\
 & 7 & $-54$ & $+20$\\
 & 9 & $-65$ & $-64$\\
 & 14 & $-140$ & .\;.\;.\;.\\
 & 15 & $-195$ & $-76$\\
\end{tabular}\end{center}

\par\indent Ist der hier sich zeigende schnell anwachsende Unterschied in AR. gegründet, so wäre ich geneigt, dies schon als eine Spur von Ellipticität der Bahn zu betrachten.“
